% =============================================================================
% EPI-AXIOME (EPISTEMOLOGIE) - Vollständige Ausformulierung
% BCM 2.2 META-Modul
% =============================================================================

\section{Epistemologie-Axiome (MA-EPI)}

Die EPI-Axiome definieren die Grenzen des Wissens, Biases und kognitive Prozesse.
Sie beantworten die Frage: \textit{Was wissen wir?}

%==============================================================================
% MA-EPI-01: Unvollständiges Wissen
%==============================================================================

\subsubsection*{MA-EPI-01: Unvollständiges Wissen}

\textbf{Statement:} Wissen über Nutzen ist immer unvollständig.

\textbf{Formel:}
\begin{equation}
K(U) \subset U
\end{equation}

\textbf{Beschreibung:}
Menschen kennen nicht alle ihre Präferenzen, nicht alle Konsequenzen ihrer Handlungen und nicht alle verfügbaren Optionen. Das subjektive Wissen $K(U)$ ist eine echte Teilmenge des objektiven Nutzens $U$. Diese Lücke ist fundamental, nicht nur praktisch.

\textbf{BCM-Relevanz:}
Epistemologisches Fundament für alle EPI-Axiome. Ohne dieses Axiom wäre vollständige Rationalität möglich. Alle kognitiven Biases sind Konsequenzen dieser fundamentalen Wissenslücke.

\textbf{BEATRIX-Relevanz:}
BEATRIX arbeitet mit geschätzten, nicht wahren Parametern:
\begin{itemize}
    \item LLM-Schätzungen sind Approximationen
    \item Unsicherheitsintervalle müssen kommuniziert werden
    \item Persona-Beschreibungen sind unvollständig
\end{itemize}

\textbf{Konsequenz bei Fehlen:}
Perfektes Wissen würde perfekte Entscheidungen ermöglichen. Informationsinterventionen wären trivial. Die gesamte Verhaltensökonomie wäre irrelevant.

\textbf{Validiert:} AWX

\textbf{Abhängigkeiten:} MA-ONT-03 (Nutzen existiert)

\textbf{Literatur:} Knight (1921); Hayek (1945)

%==============================================================================
% MA-EPI-02: Bereits in axiom_beispiele.tex
%==============================================================================

%==============================================================================
% MA-EPI-03: Irreduzible Unsicherheit
%==============================================================================

\subsubsection*{MA-EPI-03: Irreduzible Unsicherheit}

\textbf{Statement:} Es gibt fundamentale, nicht eliminierbare Unsicherheit.

\textbf{Formel:}
\begin{equation}
\exists \varepsilon: P(\varepsilon) > 0 \quad \forall \text{ Modelle}
\end{equation}

\textbf{Beschreibung:}
Selbst mit unbegrenzten Ressourcen bleibt Unsicherheit:
\begin{itemize}
    \item Ontologische Unsicherheit: Die Zukunft ist nicht determiniert
    \item Epistemische Unsicherheit: Grenzen der Messbarkeit
    \item Modell-Unsicherheit: Jedes Modell ist vereinfachend
\end{itemize}
Diese Unsicherheit ist nicht ``noch nicht gelöst'', sondern prinzipiell irreduzibel.

\textbf{BCM-Relevanz:}
Begründet Demut gegenüber Vorhersagen. Das BCM kann Wahrscheinlichkeiten angeben, nie Gewissheit. Übermäßiges Vertrauen in Modelle wird strukturell verhindert.

\textbf{BEATRIX-Relevanz:}
BEATRIX-Outputs sind probabilistisch:
\begin{itemize}
    \item Keine deterministischen Vorhersagen
    \item Konfidenzintervalle immer angeben
    \item ``Unentscheidbar'' als legitimes Ergebnis
\end{itemize}

\textbf{Konsequenz bei Fehlen:}
Hybris der Vorhersagbarkeit. Modelle würden als deterministische Wahrheiten behandelt. Risiken würden unterschätzt.

\textbf{Validiert:} META

\textbf{Abhängigkeiten:} -- (fundamentales epistemisches Axiom)

\textbf{Literatur:} Knight (1921); Taleb (2007)

%==============================================================================
% MA-EPI-04: Adaptive Heuristiken
%==============================================================================

\subsubsection*{MA-EPI-04: Adaptive Heuristiken}

\textbf{Statement:} Heuristiken können ökologisch rational sein.

\textbf{Formel:}
\begin{equation}
E[U|\text{Heuristik}] \geq E[U|\text{Optimierung}] \quad \text{in } \rho
\end{equation}

\textbf{Beschreibung:}
Heuristiken sind nicht per se Fehler. In bestimmten Umgebungen (hohe Unsicherheit, Zeitdruck) können einfache Regeln komplexe Algorithmen übertreffen. ``Fast and frugal'' Heuristiken nutzen die Struktur der Umwelt aus.

\textbf{BCM-Relevanz:}
Differenziert zwischen adaptiven und maladaptiven Abweichungen von Rationalität. Nicht jede Vereinfachung ist ein Bias. BOOST-Interventionen stärken adaptive Heuristiken.

\textbf{BEATRIX-Relevanz:}
BEATRIX unterscheidet:
\begin{itemize}
    \item Adaptive Heuristik → Nicht korrigieren
    \item Maladaptiver Bias → Intervention möglich
    \item Kontextabhängigkeit beachten
\end{itemize}

\textbf{Konsequenz bei Fehlen:}
Alle Heuristiken wären Fehler. Debiasing wäre immer angezeigt. Die ökologische Intelligenz einfacher Regeln würde übersehen.

\textbf{Validiert:} INU

\textbf{Abhängigkeiten:} MA-EPI-02

\textbf{Literatur:} Gigerenzer et al. (1999); Todd \& Gigerenzer (2012)

%==============================================================================
% MA-EPI-05: Zwei Systeme
%==============================================================================

\subsubsection*{MA-EPI-05: Zwei-Systeme-Kognition}

\textbf{Statement:} Kognition operiert in zwei Modi: System 1 (schnell, automatisch) und System 2 (langsam, deliberativ).

\textbf{Formel:}
\begin{equation}
K = K_{S1} \cup K_{S2}
\end{equation}

\textbf{Beschreibung:}
System 1: Intuitiv, assoziativ, parallelverarbeitend, mühelos, emotional getrieben.
System 2: Analytisch, regelbasiert, seriell, anstrengend, rational kontrolliert.
Die meisten Alltagsentscheidungen sind System-1-getrieben. System 2 kann aktiviert werden, ist aber ressourcenbeschränkt.

\textbf{BCM-Relevanz:}
Fundamentale Unterscheidung für Interventionsdesign:
\begin{itemize}
    \item NUDGE: Wirkt auf System 1
    \item BOOST: Stärkt System 1 oder System 2
    \item THINK: Aktiviert System 2
\end{itemize}

\textbf{BEATRIX-Relevanz:}
BEATRIX bewertet Entscheidungsmodus:
\begin{itemize}
    \item Welches System dominiert bei dieser Persona/Situation?
    \item Interventionen auf richtigen Modus abstimmen
    \item Cognitive Load beachten
\end{itemize}

\textbf{Konsequenz bei Fehlen:}
Einheitliches Kognitionsmodell. Unterschied zwischen automatischen und deliberativen Entscheidungen nicht modellierbar. Interventionstypen nicht differenzierbar.

\textbf{Validiert:} AWX, WAX

\textbf{Abhängigkeiten:} MA-EPI-02

\textbf{Literatur:} Kahneman (2011); Stanovich \& West (2000)

%==============================================================================
% MA-EPI-06: Systematische Biases
%==============================================================================

\subsubsection*{MA-EPI-06: Systematische Biases}

\textbf{Statement:} Kognitive Verzerrungen sind systematisch und vorhersagbar.

\textbf{Formel:}
\begin{equation}
E[\text{Bias}] \neq 0
\end{equation}

\textbf{Beschreibung:}
Biases sind nicht zufällige Fehler, sondern haben Richtung:
\begin{itemize}
    \item Überschätzung (Overconfidence)
    \item Unterschätzung (Neglect)
    \item Systematische Verzerrung (Framing)
\end{itemize}
Weil sie systematisch sind, sind sie vorhersagbar und potenziell korrigierbar.

\textbf{BCM-Relevanz:}
Fundament für alle spezifischen Bias-Axiome (EPI-12 bis EPI-37). Biases sind keine ``Fehler'', sondern Eigenschaften des kognitiven Systems. Das BCM integriert sie als Parameter, nicht als Störterme.

\textbf{BEATRIX-Relevanz:}
BEATRIX modelliert Biases explizit:
\begin{itemize}
    \item Bias-Anfälligkeit als Persona-Parameter
    \item Kontextuelle Bias-Aktivierung
    \item Debiasing-Potenzial bewerten
\end{itemize}

\textbf{Konsequenz bei Fehlen:}
Biases wären zufällig und nicht modellierbar. Nudges und Debiasing-Strategien hätten keine theoretische Grundlage.

\textbf{Validiert:} INU, AWX

\textbf{Abhängigkeiten:} MA-EPI-05

\textbf{Literatur:} Kahneman \& Tversky (1974); Ariely (2008)

%==============================================================================
% MA-EPI-07: AWX (Awareness)
%==============================================================================

\subsubsection*{MA-EPI-07: AWX (Awareness)}

\textbf{Statement:} Bewusstsein über Optionen und Konsequenzen variiert.

\textbf{Formel:}
\begin{equation}
AWX = f(K_V, K_C)
\end{equation}

\textbf{Beschreibung:}
AWX (Awareness) ist die erste Stufe im Verhaltenstrichter:
\begin{itemize}
    \item $K_V$: Wissen über Verhaltensoptionen (Was kann ich tun?)
    \item $K_C$: Wissen über Konsequenzen (Was passiert dann?)
\end{itemize}
Ohne Awareness keine Entscheidung. Viele Verhaltensdefizite sind Awareness-Defizite.

\textbf{BCM-Relevanz:}
Erste Variable im AWX → WAX → WTX → V Trichter. Interventionen können auf jeder Stufe ansetzen. Information und Bildung erhöhen AWX.

\textbf{BEATRIX-Relevanz:}
BEATRIX schätzt AWX-Level:
\begin{itemize}
    \item Wie informiert ist die Persona?
    \item Welche Optionen sind ihr bekannt?
    \item Werden Konsequenzen korrekt antizipiert?
\end{itemize}

\textbf{Konsequenz bei Fehlen:}
Kein Unterschied zwischen ``nicht wissen'' und ``nicht wollen''. Informationsinterventionen wären nicht gezielt einsetzbar.

\textbf{Validiert:} AWX

\textbf{Abhängigkeiten:} MA-EPI-01

\textbf{Literatur:} Prochaska \& DiClemente (1983)

%==============================================================================
% MA-EPI-08: WAX (Willingness)
%==============================================================================

\subsubsection*{MA-EPI-08: WAX (Willingness)}

\textbf{Statement:} Bereitschaft zu handeln ist eine eigenständige Variable.

\textbf{Formel:}
\begin{equation}
WAX = f(AWX, U, \Gamma)
\end{equation}

\textbf{Beschreibung:}
WAX (Willingness to Act) hängt ab von:
\begin{itemize}
    \item AWX: Muss wissen, dass Option existiert
    \item U: Muss Nutzen sehen
    \item $\Gamma$: Constraints müssen erfüllt sein
\end{itemize}
WAX ist Verhaltensintention, noch nicht Verhalten selbst.

\textbf{BCM-Relevanz:}
Zweite Stufe im Trichter. WAX kann hoch sein ohne Verhalten (Intention-Action Gap). Motivationsinterventionen erhöhen WAX.

\textbf{BEATRIX-Relevanz:}
BEATRIX bewertet WAX:
\begin{itemize}
    \item Ist die Person motiviert?
    \item Welche Barrieren senken WAX?
    \item Interventionen zur WAX-Steigerung identifizieren
\end{itemize}

\textbf{Konsequenz bei Fehlen:}
Motivation und Verhalten würden nicht unterschieden. Die wichtige Lücke zwischen ``wollen'' und ``tun'' wäre unsichtbar.

\textbf{Validiert:} WAX

\textbf{Abhängigkeiten:} MA-EPI-07, MA-ONT-03

\textbf{Literatur:} Ajzen (1991); Schwarzer (2008)

%==============================================================================
% MA-EPI-09: WTX (Transition)
%==============================================================================

\subsubsection*{MA-EPI-09: WTX (Willingness to Transition)}

\textbf{Statement:} Der Übergang von Absicht zu Verhalten ist probabilistisch.

\textbf{Formel:}
\begin{equation}
P(V|WAX) = WTX
\end{equation}

\textbf{Beschreibung:}
WTX modelliert die ``letzte Meile'': Von der Intention zur Handlung. Selbst bei hohem WAX kann WTX niedrig sein:
\begin{itemize}
    \item Prokrastination
    \item Selbstkontrollversagen
    \item Situative Barrieren
\end{itemize}

\textbf{BCM-Relevanz:}
Dritte Stufe im Trichter. WTX ist der kritische Übergangspunkt. Commitment Devices (MA-DYN-09) und Implementation Intentions (MA-DYN-18) erhöhen WTX.

\textbf{BEATRIX-Relevanz:}
BEATRIX schätzt WTX:
\begin{itemize}
    \item Wie wahrscheinlich ist die Umsetzung?
    \item Welche Transitonsbarrieren existieren?
    \item Braucht es einen ``letzten Schubs''?
\end{itemize}

\textbf{Konsequenz bei Fehlen:}
Intention würde mit Verhalten gleichgesetzt. Die Intention-Action Gap (MA-EPI-11) wäre nicht modellierbar.

\textbf{Validiert:} WTX

\textbf{Abhängigkeiten:} MA-EPI-08

\textbf{Literatur:} Sheeran (2002); Gollwitzer \& Sheeran (2006)

%==============================================================================
% MA-EPI-10: AWX-WAX-WTX Sequenz
%==============================================================================

\subsubsection*{MA-EPI-10: Verhaltenssequenz AWX → WAX → WTX → V}

\textbf{Statement:} Verhalten folgt einer sequenziellen Kette.

\textbf{Formel:}
\begin{equation}
V = WTX(WAX(AWX(K)))
\end{equation}

\textbf{Beschreibung:}
Die Kette ist kausal geordnet:
\begin{enumerate}
    \item Ohne AWX kein WAX (Was ich nicht kenne, kann ich nicht wollen)
    \item Ohne WAX kein WTX (Ohne Motivation keine Handlung)
    \item Ohne WTX kein V (Ohne Umsetzung kein Verhalten)
\end{enumerate}
Jede Stufe ist notwendig, keine ist hinreichend.

\textbf{BCM-Relevanz:}
Ermöglicht Diagnose: Wo liegt das Problem? Interventionen müssen auf die richtige Stufe zielen:
\begin{itemize}
    \item AWX-Defizit → Information
    \item WAX-Defizit → Motivation
    \item WTX-Defizit → Facilitation
\end{itemize}

\textbf{BEATRIX-Relevanz:}
BEATRIX identifiziert den Engpass:
\begin{itemize}
    \item Welche Stufe ist der Bottleneck?
    \item Interventionen entsprechend priorisieren
    \item Kaskadeneffekte modellieren
\end{itemize}

\textbf{Konsequenz bei Fehlen:}
Alle Verhaltensdefizite würden gleich behandelt. Intervention ohne Diagnose.

\textbf{Validiert:} AWX, WAX, WTX

\textbf{Abhängigkeiten:} MA-EPI-07, MA-EPI-08, MA-EPI-09

\textbf{Literatur:} Michie et al. (2011)

%==============================================================================
% MA-EPI-11: Intention-Action Gap
%==============================================================================

\subsubsection*{MA-EPI-11: Intention-Action Gap}

\textbf{Statement:} Absicht führt nicht immer zu Verhalten.

\textbf{Formel:}
\begin{equation}
P(V|WAX=1) < 1
\end{equation}

\textbf{Beschreibung:}
Selbst bei maximaler Motivation wird das Verhalten nicht immer umgesetzt:
\begin{itemize}
    \item 47\% der Intentions werden nicht umgesetzt (Sheeran, 2002)
    \item Gründe: Vergessen, Aufschub, situative Hindernisse
    \item Chronische Gap bei Gesundheits- und Umweltverhalten
\end{itemize}

\textbf{BCM-Relevanz:}
Zentral für WTX-Forschung. Die Gap ist nicht Motivationsdefizit, sondern Ausführungsdefizit. Andere Interventionen nötig: Commitment, Defaults, Erinnerungen.

\textbf{BEATRIX-Relevanz:}
BEATRIX quantifiziert die Gap:
\begin{itemize}
    \item Wie groß ist die typische Gap bei dieser Persona/Verhalten?
    \item Welche Interventionen schließen die Gap?
    \item Modelliert als $P(V|WAX)$
\end{itemize}

\textbf{Konsequenz bei Fehlen:}
$WAX = 1$ würde $V = 1$ implizieren. Umsetzungsprobleme wären unsichtbar.

\textbf{Validiert:} WTX

\textbf{Abhängigkeiten:} MA-EPI-09

\textbf{Literatur:} Sheeran (2002); Webb \& Sheeran (2006)

%==============================================================================
% MA-EPI-12: Verfügbarkeitsheuristik
%==============================================================================

\subsubsection*{MA-EPI-12: Verfügbarkeitsheuristik}

\textbf{Statement:} Leicht Abrufbares wird als wahrscheinlicher eingeschätzt.

\textbf{Formel:}
\begin{equation}
P_{\text{subj}}(E) = f(\text{Accessibility}(E))
\end{equation}

\textbf{Beschreibung:}
Die Verfügbarkeitsheuristik ersetzt ``Wie wahrscheinlich?'' durch ``Wie leicht fällt mir ein Beispiel ein?'':
\begin{itemize}
    \item Medial präsente Risiken (Terror) überschätzt
    \item Alltägliche Risiken (Autounfall) unterschätzt
    \item Persönliche Erfahrungen übergewichtet
\end{itemize}

\textbf{BCM-Relevanz:}
Erklärt systematische Risikofehleinschätzungen. Kann für Interventionen genutzt werden: Saliente Beispiele erhöhen wahrgenommene Wahrscheinlichkeit.

\textbf{BEATRIX-Relevanz:}
BEATRIX berücksichtigt Verfügbarkeit:
\begin{itemize}
    \item Welche Beispiele sind für die Persona salient?
    \item Welche Medienexposition?
    \item Können saliente Counter-Examples präsentiert werden?
\end{itemize}

\textbf{Konsequenz bei Fehlen:}
Wahrscheinlichkeitsurteile wären Bayes-rational. Media-Effekte auf Risikowahrnehmung unerklärlich.

\textbf{Validiert:} AWX

\textbf{Abhängigkeiten:} MA-EPI-05

\textbf{Literatur:} Tversky \& Kahneman (1973)

%==============================================================================
% MA-EPI-13: Framing-Effekte
%==============================================================================

\subsubsection*{MA-EPI-13: Framing-Effekte}

\textbf{Statement:} Die Darstellung beeinflusst die Entscheidung.

\textbf{Formel:}
\begin{equation}
U(V|\text{Frame}_A) \neq U(V|\text{Frame}_B)
\end{equation}

\textbf{Beschreibung:}
Dieselbe Information, unterschiedlich präsentiert, führt zu unterschiedlichen Entscheidungen:
\begin{itemize}
    \item ``90\% Überlebensrate'' vs. ``10\% Sterberate''
    \item ``500€ Rabatt'' vs. ``500€ weniger zahlen''
    \item Gewinn-Frame vs. Verlust-Frame
\end{itemize}

\textbf{BCM-Relevanz:}
Mächtiges Interventionsinstrument. NUDGE-Interventionen nutzen häufig Framing. Ethisch: Transparenz über Frame-Wahl nötig.

\textbf{BEATRIX-Relevanz:}
BEATRIX analysiert Frame-Sensitivität:
\begin{itemize}
    \item Wie Frame-sensitiv ist die Persona?
    \item Welcher Frame ist für das Zielverhalten optimal?
    \item Ethische Bewertung des Framings
\end{itemize}

\textbf{Konsequenz bei Fehlen:}
Präsentation wäre irrelevant. Nur Inhalt zählt (homo oeconomicus). Kommunikationsdesign hätte keinen Wert.

\textbf{Validiert:} INU, AWX

\textbf{Abhängigkeiten:} MA-EPI-06

\textbf{Literatur:} Tversky \& Kahneman (1981)

%==============================================================================
% MA-EPI-14: Anchoring
%==============================================================================

\subsubsection*{MA-EPI-14: Anchoring}

\textbf{Statement:} Anker beeinflussen numerische Schätzungen.

\textbf{Formel:}
\begin{equation}
E[X|\text{Anchor}] = f(\text{Anchor})
\end{equation}

\textbf{Beschreibung:}
Beliebige Zahlen (auch irrelevante) ziehen Schätzungen in ihre Richtung:
\begin{itemize}
    \item Preisverhandlungen: Erstes Angebot dominiert
    \item Spendenaufrufe: Vorgeschlagene Beträge wirken als Anker
    \item Selbst offensichtlich irrelevante Zahlen haben Wirkung
\end{itemize}

\textbf{BCM-Relevanz:}
Wird in NUDGE-Interventionen genutzt. Default-Werte sind Anker. Ethik: Anker sollten nicht manipulativ sein.

\textbf{BEATRIX-Relevanz:}
BEATRIX berücksichtigt Ankereffekte:
\begin{itemize}
    \item Welche Anker sind in der Situation präsent?
    \item Wie kann Anchoring für positive Verhaltensänderung genutzt werden?
    \item Debiasing durch multiple Anker
\end{itemize}

\textbf{Konsequenz bei Fehlen:}
Schätzungen wären unbeeinflussbar durch Kontext. Defaults und Vorschläge hätten keine Wirkung.

\textbf{Validiert:} INU

\textbf{Abhängigkeiten:} MA-EPI-06

\textbf{Literatur:} Tversky \& Kahneman (1974); Ariely et al. (2003)

%==============================================================================
% MA-EPI-15: Bereits in axiom_beispiele.tex
%==============================================================================

%==============================================================================
% MA-EPI-16: Present Bias
%==============================================================================

\subsubsection*{MA-EPI-16: Present Bias}

\textbf{Statement:} Die Gegenwart wird übergewichtet.

\textbf{Formel:}
\begin{equation}
\delta_{\text{now}} > \delta_{\text{future}}
\end{equation}

\textbf{Beschreibung:}
Unmittelbare Belohnungen werden unverhältnismäßig stark gewichtet:
\begin{itemize}
    \item ``Jetzt 50€'' wird ``in einer Woche 55€'' vorgezogen
    \item Aber: ``In einem Jahr 50€'' vs. ``in einem Jahr + einer Woche 55€'' -- hier wird gewartet
    \item Dies ist zeitinkonsistent (dynamische Inkonsistenz)
\end{itemize}

\textbf{BCM-Relevanz:}
Erklärt Prokrastination, Suchtverhalten und Spar-Defizite. Present Bias ist der Kern von MA-DYN-14 (hyperbolisches Discounting).

\textbf{BEATRIX-Relevanz:}
BEATRIX schätzt Present-Bias-Stärke:
\begin{itemize}
    \item Hinweise auf Impulsivität in Persona
    \item Interventionen mit sofortiger vs. verzögerter Belohnung
    \item Commitment Devices für Hochimpulsive
\end{itemize}

\textbf{Konsequenz bei Fehlen:}
Exponentielles Discounting (zeitkonsistent). Prokrastination wäre unerklärlich. Commitment Devices nutzlos.

\textbf{Validiert:} INU, WTX

\textbf{Abhängigkeiten:} MA-ONT-05

\textbf{Literatur:} Laibson (1997); O'Donoghue \& Rabin (1999)

%==============================================================================
% MA-EPI-17: Status Quo Bias
%==============================================================================

\subsubsection*{MA-EPI-17: Status Quo Bias}

\textbf{Statement:} Der aktuelle Zustand wird bevorzugt.

\textbf{Formel:}
\begin{equation}
U(\text{Status Quo}) + \text{Bias} > U(\text{Alternative})
\end{equation}

\textbf{Beschreibung:}
Menschen bleiben bei dem, was sie haben, selbst wenn Alternativen objektiv besser sind:
\begin{itemize}
    \item Nichtwechseln von Stromanbietern
    \item Beibehalten suboptimaler Voreinstellungen
    \item Festhalten an Investitionsentscheidungen
\end{itemize}

\textbf{BCM-Relevanz:}
Erklärt die Macht von Defaults. Wenn der Default ``Organspende'' ist, spenden mehr. Status Quo Bias + Loss Aversion verstärken sich.

\textbf{BEATRIX-Relevanz:}
BEATRIX nutzt Status Quo Bias:
\begin{itemize}
    \item Was ist der aktuelle Default?
    \item Kann der Default geändert werden?
    \item Bei SQ-Bias-anfälliger Persona: Opt-Out statt Opt-In
\end{itemize}

\textbf{Konsequenz bei Fehlen:}
Defaults wären irrelevant. Alle Optionen würden neutral bewertet.

\textbf{Validiert:} INU, WAX

\textbf{Abhängigkeiten:} MA-EPI-15

\textbf{Literatur:} Samuelson \& Zeckhauser (1988); Kahneman et al. (1991)

%==============================================================================
% MA-EPI-18: Overconfidence
%==============================================================================

\subsubsection*{MA-EPI-18: Overconfidence}

\textbf{Statement:} Eigene Fähigkeiten werden überschätzt.

\textbf{Formel:}
\begin{equation}
P_{\text{subj}}(\text{Success}) > P_{\text{obj}}(\text{Success})
\end{equation}

\textbf{Beschreibung:}
Die meisten Menschen glauben, überdurchschnittlich zu sein:
\begin{itemize}
    \item 93\% der US-Fahrer halten sich für ``überdurchschnittlich''
    \item Überschätzung eigener Prognosegenauigkeit
    \item ``Planning Fallacy'': Projekte dauern länger als geplant
\end{itemize}

\textbf{BCM-Relevanz:}
Overconfidence beeinflusst Risikobereitschaft und Planungsverhalten. Kann zu suboptimalen Entscheidungen führen (Unternehmensgründungen, Investitionen).

\textbf{BEATRIX-Relevanz:}
BEATRIX korrigiert für Overconfidence:
\begin{itemize}
    \item Persona mit hohem Overconfidence → Risikowarnung
    \item Zeitschätzungen nach oben korrigieren
    \item Objektive Benchmarks präsentieren
\end{itemize}

\textbf{Konsequenz bei Fehlen:}
Selbsteinschätzungen würden als akkurat behandelt. Unrealistische Pläne blieben unkorrigiert.

\textbf{Validiert:} INU, AWX

\textbf{Abhängigkeiten:} MA-EPI-06

\textbf{Literatur:} Moore \& Healy (2008); Buehler et al. (1994)

%==============================================================================
% MA-EPI-19: Lernen ist möglich
%==============================================================================

\subsubsection*{MA-EPI-19: Lernen ist möglich}

\textbf{Statement:} Wissen kann akkumuliert werden.

\textbf{Formel:}
\begin{equation}
K(\tau + \Delta\tau) \supseteq K(\tau)
\end{equation}

\textbf{Beschreibung:}
Trotz fundamentaler Wissensgrenzen (MA-EPI-01) kann Wissen erweitert werden:
\begin{itemize}
    \item Erfahrungslernen
    \item Instruktion und Bildung
    \item Feedback-basiertes Lernen
\end{itemize}
Die Wissensmenge ist monoton nicht-abnehmend (Vergessen ignoriert).

\textbf{BCM-Relevanz:}
Begründet BOOST- und THINK-Interventionen. Wenn Lernen unmöglich wäre, könnten nur Constraints und Nudges wirken.

\textbf{BEATRIX-Relevanz:}
BEATRIX modelliert Lernpotenzial:
\begin{itemize}
    \item Wie lernfähig ist die Persona?
    \item Welche Lerninterventionen sind effektiv?
    \item Erwartete Wissensakkumulation über Zeit
\end{itemize}

\textbf{Konsequenz bei Fehlen:}
Statisches Wissen. Bildungsinterventionen wären wirkungslos. Keine Kompetenzentwicklung modellierbar.

\textbf{Validiert:} AWX

\textbf{Abhängigkeiten:} MA-EPI-01, MA-ONT-05

\textbf{Literatur:} Anderson (1982); Ericsson \& Charness (1994)

%==============================================================================
% MA-EPI-20: Feedback verbessert Schätzungen
%==============================================================================

\subsubsection*{MA-EPI-20: Feedback verbessert Schätzungen}

\textbf{Statement:} Mit Feedback werden Prognosen akkurater.

\textbf{Formel:}
\begin{equation}
|K(\tau+1) - U| < |K(\tau) - U| \quad \text{mit Feedback}
\end{equation}

\textbf{Beschreibung:}
Feedback ermöglicht Kalibrierung:
\begin{itemize}
    \item Überprüfung von Vorhersagen an der Realität
    \item Korrektur systematischer Biases
    \item Verbesserung der Metakognition
\end{itemize}
Voraussetzung: Feedback muss zeitnah, spezifisch und glaubwürdig sein.

\textbf{BCM-Relevanz:}
Begründet Feedback-Interventionen. Real-time Feedback (Stromverbrauch, Fitness) kann Verhalten ändern.

\textbf{BEATRIX-Relevanz:}
BEATRIX empfiehlt Feedback-Mechanismen:
\begin{itemize}
    \item Welches Feedback wäre hilfreich?
    \item Wie schnell und spezifisch sollte es sein?
    \item Feedback-Loop in Interventionsdesign integrieren
\end{itemize}

\textbf{Konsequenz bei Fehlen:}
Feedback wäre wirkungslos. Keine Lernkurven. Statische Biases.

\textbf{Validiert:} AWX

\textbf{Abhängigkeiten:} MA-EPI-19

\textbf{Literatur:} Kluger \& DeNisi (1996); Hattie \& Timperley (2007)

%==============================================================================
% MA-EPI-21: Gewohnheiten reduzieren Kognition
%==============================================================================

\subsubsection*{MA-EPI-21: Gewohnheiten reduzieren kognitive Last}

\textbf{Statement:} Automatisierung spart kognitive Ressourcen.

\textbf{Formel:}
\begin{equation}
K_{\text{needed}}(V_{\text{habit}}) < K_{\text{needed}}(V_{\text{new}})
\end{equation}

\textbf{Beschreibung:}
Gewohnheiten sind kognitiv effizient:
\begin{itemize}
    \item Automatische Ausführung ohne bewusste Entscheidung
    \item Freisetzung kognitiver Ressourcen für andere Aufgaben
    \item Stabilität auch unter Stress
\end{itemize}
Nachteile: Schwer zu ändern, können maladaptiv sein.

\textbf{BCM-Relevanz:}
Erklärt Persistenz von Verhalten. Gewohnheitsänderung erfordert temporär erhöhte kognitive Ressourcen. DESIGN-Interventionen können Habit-Cues unterbrechen.

\textbf{BEATRIX-Relevanz:}
BEATRIX analysiert Gewohnheitsgrad:
\begin{itemize}
    \item Ist das Verhalten habituiert?
    \item Welche Cues triggern die Gewohnheit?
    \item Interventionsstrategie: Habit Disruption oder Habit Formation
\end{itemize}

\textbf{Konsequenz bei Fehlen:}
Alle Entscheidungen würden deliberativ behandelt. Automatisches Verhalten nicht modellierbar.

\textbf{Validiert:} WTX

\textbf{Abhängigkeiten:} MA-EPI-05

\textbf{Literatur:} Wood \& Neal (2007); Verplanken \& Wood (2006)

%==============================================================================
% MA-EPI-22: Debiasing ist schwierig
%==============================================================================

\subsubsection*{MA-EPI-22: Debiasing ist schwierig}

\textbf{Statement:} Kognitive Biases sind persistent.

\textbf{Formel:}
\begin{equation}
\frac{\partial \text{Bias}}{\partial \text{Intervention}} < \text{expected}
\end{equation}

\textbf{Beschreibung:}
Biases lassen sich nicht einfach ``wegtrainieren'':
\begin{itemize}
    \item Wissen über einen Bias schützt nicht davor
    \item Effekte von Debiasing-Training sind kurzlebig
    \item Biases kehren unter Stress zurück
\end{itemize}
Besser: Umgebungen designen, die Biases irrelevant machen.

\textbf{BCM-Relevanz:}
Erklärt Grenzen von THINK-Interventionen. Reine Aufklärung reicht nicht. DESIGN-Interventionen, die Bias-Effekte neutralisieren, sind oft wirksamer.

\textbf{BEATRIX-Relevanz:}
BEATRIX warnt vor Überoptimismus bezüglich Debiasing:
\begin{itemize}
    \item Debiasing-Interventionen mit Vorsicht empfehlen
    \item Alternative: Choice Architecture
    \item Langfristige Wirkung skeptisch bewerten
\end{itemize}

\textbf{Konsequenz bei Fehlen:}
Biases würden als leicht korrigierbar angenommen. Überinvestition in Aufklärung statt Umgebungsdesign.

\textbf{Validiert:} AWX

\textbf{Abhängigkeiten:} MA-EPI-06

\textbf{Literatur:} Fischhoff (1982); Lilienfeld et al. (2009)

%==============================================================================
% MA-EPI-23: Exponentielle Sättigung
%==============================================================================

\subsubsection*{MA-EPI-23: Exponentielle Sättigung (MA\_0.ES)}

\textbf{Statement:} Systeme mit proportionaler Unsicherheitsreduktion konvergieren exponentiell.

\textbf{Formel:}
\begin{equation}
\frac{dU}{dt} = -k \cdot U \quad \Rightarrow \quad U(t) = U_0 \cdot e^{-kt}
\end{equation}

\textbf{Beschreibung:}
Wenn Unsicherheit proportional zu sich selbst reduziert wird:
\begin{itemize}
    \item Schneller Anfangsprogress
    \item Asymptotische Annäherung an Null
    \item Übergangspunkt bei ~63\% Abbau (Rest $\approx 1/e$)
\end{itemize}
Anwendung: Lernkurven, Informationsakkumulation, Habituation.

\textbf{BCM-Relevanz:}
Mathematische Grundlage für viele dynamische Prozesse im BCM. Erklärt, warum initiale Interventionseffekte groß sind, dann abnehmen.

\textbf{BEATRIX-Relevanz:}
BEATRIX modelliert Sättigungseffekte:
\begin{itemize}
    \item Wann ist der Übergangspunkt erreicht?
    \item Ab wann lohnen weitere Interventionen nicht mehr?
    \item Optimaler Interventionszeitpunkt
\end{itemize}

\textbf{Konsequenz bei Fehlen:}
Lineare Extrapolation statt Sättigung. Langfristprognosen wären falsch.

\textbf{Validiert:} META

\textbf{Abhängigkeiten:} MA-EPI-19

\textbf{Literatur:} Eigenes Konstrukt; vgl. Newell \& Rosenbloom (1981)

%==============================================================================
% MA-EPI-24: Bayes'sches Updating
%==============================================================================

\subsubsection*{MA-EPI-24: Bayes'sches Updating}

\textbf{Statement:} Rationale Belief-Updates folgen dem Bayes-Theorem.

\textbf{Formel:}
\begin{equation}
P(H|E) = \frac{P(E|H) \cdot P(H)}{P(E)}
\end{equation}

\textbf{Beschreibung:}
Bayes beschreibt, wie Überzeugungen angesichts neuer Evidenz aktualisiert werden sollten:
\begin{itemize}
    \item Prior $P(H)$: Vorherige Überzeugung
    \item Likelihood $P(E|H)$: Wie wahrscheinlich ist die Evidenz unter der Hypothese?
    \item Posterior $P(H|E)$: Aktualisierte Überzeugung
\end{itemize}

\textbf{BCM-Relevanz:}
Normativer Benchmark für Lernprozesse. Abweichungen von Bayes (Base Rate Neglect, Confirmation Bias) sind als Biases modelliert.

\textbf{BEATRIX-Relevanz:}
BEATRIX verwendet Bayes für:
\begin{itemize}
    \item Eigene Inferenz über Persona-Parameter
    \item Modellierung, wie die Persona lernt
    \item Benchmark für Bias-Identifikation
\end{itemize}

\textbf{Konsequenz bei Fehlen:}
Kein normativer Standard für Lernen. Biases wären nicht als Abweichungen identifizierbar.

\textbf{Validiert:} AWX

\textbf{Abhängigkeiten:} MA-EPI-19

\textbf{Literatur:} Bayes (1763); Edwards et al. (1963)

%==============================================================================
% MA-EPI-25: Confirmation Bias
%==============================================================================

\subsubsection*{MA-EPI-25: Confirmation Bias}

\textbf{Statement:} Bestätigende Information wird bevorzugt gesucht und verarbeitet.

\textbf{Formel:}
\begin{equation}
P(\text{Search}|\text{Confirms}) > P(\text{Search}|\text{Disconfirms})
\end{equation}

\textbf{Beschreibung:}
Menschen suchen aktiv nach Information, die bestehende Überzeugungen bestätigt:
\begin{itemize}
    \item Selektive Mediennutzung (Echo Chambers)
    \item Asymmetrische Evidenzbewertung
    \item Immunisierung gegen Gegenargumente
\end{itemize}

\textbf{BCM-Relevanz:}
Erklärt Persistenz von Fehlüberzeugungen. Information allein ändert Beliefs nicht. Dialogische Interventionen nötig.

\textbf{BEATRIX-Relevanz:}
BEATRIX berücksichtigt Confirmation Bias:
\begin{itemize}
    \item Welche Überzeugungen hat die Persona?
    \item Wird neue Information abgelehnt werden?
    \item Framing so, dass es nicht als Angriff wirkt
\end{itemize}

\textbf{Konsequenz bei Fehlen:}
Information würde neutral verarbeitet. Fake News und Polarisierung unerklärlich.

\textbf{Validiert:} AWX

\textbf{Abhängigkeiten:} MA-EPI-06

\textbf{Literatur:} Nickerson (1998); Klayman (1995)

%==============================================================================
% MA-EPI-26: Bereits in axiom_beispiele.tex
%==============================================================================

%==============================================================================
% MA-EPI-27: Representativeness Heuristic
%==============================================================================

\subsubsection*{MA-EPI-27: Representativeness Heuristic}

\textbf{Statement:} Urteile basieren auf Ähnlichkeit zu Prototypen.

\textbf{Formel:}
\begin{equation}
\hat{P}(A \in \text{Category}) = f(\text{Similarity}(A, \text{Prototype}))
\end{equation}

\textbf{Beschreibung:}
Die Repräsentativitätsheuristik ersetzt Wahrscheinlichkeit durch Ähnlichkeit:
\begin{itemize}
    \item ``Linda Problem'': Konjunktionsfehlschluss
    \item Stereotypen-basierte Urteile
    \item Ignorieren von Base Rates (MA-EPI-26)
\end{itemize}

\textbf{BCM-Relevanz:}
Fundamental für Typenklassifikation. Menschen kategorisieren andere (und sich selbst) nach Prototypen. Kann zu systematischen Fehlklassifikationen führen.

\textbf{BEATRIX-Relevanz:}
BEATRIX korrigiert für Representativeness:
\begin{itemize}
    \item Prototypen explizit definieren
    \item Base Rates gewichten
    \item Vorsicht bei ``passt perfekt''-Urteilen
\end{itemize}

\textbf{Konsequenz bei Fehlen:}
Kategorisierung wäre rein probabilistisch. Intuitive Typisierung nicht modellierbar.

\textbf{Validiert:} AWX, ESTIMATION

\textbf{Abhängigkeiten:} MA-EPI-04

\textbf{Literatur:} Kahneman \& Tversky (1972)

%==============================================================================
% MA-EPI-28: Sunk Cost Fallacy
%==============================================================================

\subsubsection*{MA-EPI-28: Sunk Cost Fallacy}

\textbf{Statement:} Versunkene Kosten beeinflussen zukünftige Entscheidungen.

\textbf{Formel:}
\begin{equation}
U(\text{Continue}|\text{SunkCost}) > U(\text{Continue}|\text{NoSunkCost})
\end{equation}

\textbf{Beschreibung:}
Rationale Entscheidungen sollten nur zukünftige Kosten und Nutzen berücksichtigen. Aber:
\begin{itemize}
    \item ``Ich habe schon so viel investiert''
    \item Weitermachen bei gescheiterten Projekten
    \item Essen aufessen ``weil bezahlt''
\end{itemize}

\textbf{BCM-Relevanz:}
Erklärt irrationale Persistenz. Relevant für Commitment und Eskalation. Kann Verhaltensänderung blockieren (``habe schon so lange geraucht'').

\textbf{BEATRIX-Relevanz:}
BEATRIX identifiziert Sunk-Cost-Fallen:
\begin{itemize}
    \item Welche vergangenen Investitionen blockieren Änderung?
    \item Reframing: Zukunftsorientierte Perspektive fördern
    \item ``Sunk Costs sind sunk'' -- Debiasing-Intervention
\end{itemize}

\textbf{Konsequenz bei Fehlen:}
Vergangene Investitionen wären irrelevant. Eskalation von Commitment unerklärlich.

\textbf{Validiert:} INU

\textbf{Abhängigkeiten:} MA-EPI-15

\textbf{Literatur:} Arkes \& Blumer (1985); Thaler (1980)

%==============================================================================
% MA-EPI-29: Endowment Effect
%==============================================================================

\subsubsection*{MA-EPI-29: Endowment Effect}

\textbf{Statement:} Eigene Güter werden höher bewertet als identische fremde.

\textbf{Formel:}
\begin{equation}
WTA > WTP, \quad \text{typisch: } WTA/WTP \approx 2
\end{equation}

\textbf{Beschreibung:}
Die Verkaufsbereitschaft (WTA = Willingness to Accept) ist systematisch höher als die Kaufbereitschaft (WTP = Willingness to Pay) für dasselbe Gut:
\begin{itemize}
    \item Kaffeebecherexperiment: WTA 2x so hoch wie WTP
    \item ``Mein Kind ist das klügste''
    \item Überbewertung eigener Ideen/Arbeit
\end{itemize}

\textbf{BCM-Relevanz:}
Konsequenz von Loss Aversion (MA-EPI-15). Abgeben ist ein Verlust. Erklärt Marktfriktionen und suboptimalen Handel.

\textbf{BEATRIX-Relevanz:}
BEATRIX berücksichtigt Endowment:
\begin{itemize}
    \item Was ``besitzt'' die Persona (materiell und immateriell)?
    \item Wie stark ist der Endowment Effect bei ihr?
    \item Interventionen: Probebesitz vs. direkter Kauf
\end{itemize}

\textbf{Konsequenz bei Fehlen:}
WTA = WTP (ökonomisches Standardmodell). Besitzeffekte unerklärlich.

\textbf{Validiert:} INU

\textbf{Abhängigkeiten:} MA-EPI-15

\textbf{Literatur:} Kahneman, Knetsch \& Thaler (1990)

%==============================================================================
% MA-EPI-30: Optimism Bias
%==============================================================================

\subsubsection*{MA-EPI-30: Optimism Bias}

\textbf{Statement:} Positive Outcomes werden für sich selbst überschätzt.

\textbf{Formel:}
\begin{equation}
\hat{P}(\text{Positive}|\text{Self}) > P(\text{Positive})
\end{equation}

\textbf{Beschreibung:}
Unrealistischer Optimismus für eigene Zukunft:
\begin{itemize}
    \item ``Mir passiert das nicht'' (Unfälle, Krankheiten)
    \item Überschätzung von Karrierechancen
    \item Unterschätzung eigener Risiken
\end{itemize}

\textbf{BCM-Relevanz:}
Erklärt Unterversicherung, Risikoignoranz und mangelnde Prävention. Optimism Bias ist adaptiv (Wohlbefinden), aber kann zu Fehlentscheidungen führen.

\textbf{BEATRIX-Relevanz:}
BEATRIX korrigiert für Optimism Bias:
\begin{itemize}
    \item Risikowahrnehmung nach oben korrigieren
    \item Personalisierte Risikostatistiken präsentieren
    \item ``Das kann mir passieren''-Interventionen
\end{itemize}

\textbf{Konsequenz bei Fehlen:}
Selbst-Risiken würden korrekt geschätzt. Präventionslücken unerklärlich.

\textbf{Validiert:} AWX, IDN

\textbf{Abhängigkeiten:} MA-EPI-06

\textbf{Literatur:} Weinstein (1980); Sharot (2011)

%==============================================================================
% MA-EPI-31: Self-Serving Bias
%==============================================================================

\subsubsection*{MA-EPI-31: Self-Serving Bias}

\textbf{Statement:} Erfolge werden internal, Misserfolge external attribuiert.

\textbf{Formel:}
\begin{equation}
P(\text{Internal}|\text{Success}) > P(\text{Internal}|\text{Failure})
\end{equation}

\textbf{Beschreibung:}
Attributionsmuster zum Selbstwertschutz:
\begin{itemize}
    \item Erfolg: ``Ich bin gut''
    \item Misserfolg: ``Pech gehabt'' oder ``unfair''
\end{itemize}
Dient der Selbstwerterhaltung, behindert aber Lernen aus Fehlern.

\textbf{BCM-Relevanz:}
Kern-Mechanismus für Motivated Inference im IDN-Modul. Identitätsschutz durch selektive Attribution. Relevant für Feedback-Interventionen.

\textbf{BEATRIX-Relevanz:}
BEATRIX berücksichtigt Self-Serving Bias:
\begin{itemize}
    \item Wie wird die Persona Misserfolge attribuieren?
    \item Feedback so gestalten, dass Selbstwert nicht bedroht wird
    \item Lernen aus Fehlern ermöglichen
\end{itemize}

\textbf{Konsequenz bei Fehlen:}
Symmetrische Attribution. Selbstwertschutz-Dynamik nicht modellierbar.

\textbf{Validiert:} IDN

\textbf{Abhängigkeiten:} MA-EPI-06

\textbf{Literatur:} Miller \& Ross (1975); Mezulis et al. (2004)

%==============================================================================
% MA-EPI-32: Social Desirability Bias
%==============================================================================

\subsubsection*{MA-EPI-32: Social Desirability Bias}

\textbf{Statement:} Selbstberichte sind Richtung sozialer Erwünschtheit verzerrt.

\textbf{Formel:}
\begin{equation}
\text{Report}(V) = V_{\text{true}} + \varepsilon_{SD}, \quad \varepsilon_{SD} = f(\text{Visibility}, \text{Norm})
\end{equation}

\textbf{Beschreibung:}
Menschen berichten nicht wahrheitsgemäß, wenn die Wahrheit sozial unerwünscht ist:
\begin{itemize}
    \item Übertreibung von prosozialem Verhalten
    \item Untertreibung von Tabus (Alkohol, Rassismus)
    \item Stärker bei beobachteten Antworten
\end{itemize}

\textbf{BCM-Relevanz:}
\textbf{Kritisch für BEATRIX!} Persona-Beschreibungen sind systematisch verzerrt. Das LLM sieht eine ``aufgehübschte'' Version.

\textbf{BEATRIX-Relevanz:}
BEATRIX korrigiert für Social Desirability:
\begin{itemize}
    \item Persona-Angaben diskontieren
    \item Verhalten > Selbstbericht gewichten
    \item ``True self'' liegt näher am Mittelwert
\end{itemize}

\textbf{Konsequenz bei Fehlen:}
Persona-Beschreibungen würden für bare Münze genommen. Systematische Überschätzung von ``guten'' Eigenschaften.

\textbf{Validiert:} AWX, ESTIMATION

\textbf{Abhängigkeiten:} MA-EPI-06, MA-ONT-11

\textbf{Literatur:} Crowne \& Marlowe (1960); Paulhus (1984)

%==============================================================================
% MA-EPI-33: Omission Bias
%==============================================================================

\subsubsection*{MA-EPI-33: Omission Bias}

\textbf{Statement:} Unterlassung wird milder bewertet als aktive Handlung.

\textbf{Formel:}
\begin{equation}
U(\text{Outcome}|\text{Omission}) > U(\text{Outcome}|\text{Commission})
\end{equation}

\textbf{Beschreibung:}
Bei gleichem negativen Outcome wird Nichtstun weniger negativ bewertet als aktives Handeln:
\begin{itemize}
    \item Trolley-Problem: Nicht-Umleiten vs. aktives Stoßen
    \item Impfverweigerung: ``Nebenwirkung durch Impfung'' > ``Krankheit durch Nicht-Impfen''
\end{itemize}

\textbf{BCM-Relevanz:}
Relevant für ethische Bewertung von Interventionen. Opt-Out wird anders bewertet als Opt-In, selbst bei gleichem Ergebnis.

\textbf{BEATRIX-Relevanz:}
BEATRIX nutzt Omission Bias:
\begin{itemize}
    \item Default-Interventionen sind ``Omission'' für den Nutzer
    \item Aktive Entscheidung fühlt sich ``riskanter'' an
    \item Framing: ``Nichtstun hat Konsequenzen''
\end{itemize}

\textbf{Konsequenz bei Fehlen:}
Handlung und Unterlassung würden gleich bewertet. Default-Effekte wären rein durch Effort erklärbar.

\textbf{Validiert:} GOV, IDN

\textbf{Abhängigkeiten:} MA-EPI-06

\textbf{Literatur:} Spranca, Minsk \& Baron (1991); Ritov \& Baron (1990)

%==============================================================================
% MA-EPI-34: Moral Licensing
%==============================================================================

\subsubsection*{MA-EPI-34: Moral Licensing}

\textbf{Statement:} Gute Taten lizenzieren nachfolgende schlechte.

\textbf{Formel:}
\begin{equation}
P(\text{BadAction}|\text{PriorGoodAction}) > P(\text{BadAction}|\text{NoPriorAction})
\end{equation}

\textbf{Beschreibung:}
Moral Licensing ist ein negativer Spillover-Effekt:
\begin{itemize}
    \item ``Ich habe Sport gemacht, also kann ich jetzt Kuchen essen''
    \item ``Ich habe gespendet, also muss ich nicht recyceln''
    \item Selbstbild-Kredit wird aufgebaut und ausgegeben
\end{itemize}

\textbf{BCM-Relevanz:}
Identitäts-Dynamik: Gute Taten stärken das Selbstbild als ``guter Mensch'', was temporär schlechtes Verhalten erlaubt. Wichtig für Interventionskombinationen.

\textbf{BEATRIX-Relevanz:}
BEATRIX warnt vor Moral Licensing:
\begin{itemize}
    \item Nach erfolgreicher Intervention 1: Risiko für Backlash in Bereich 2
    \item Nicht mehrere ``gute'' Verhaltensweisen simultan einfordern
    \item Identitätsbasierte Interventionen können Licensing auslösen
\end{itemize}

\textbf{Konsequenz bei Fehlen:}
Positiver Spillover würde überschätzt. Interventionspakete könnten kontraproduktiv sein.

\textbf{Validiert:} IDN, KNU

\textbf{Abhängigkeiten:} MA-EPI-06

\textbf{Literatur:} Merritt, Effron \& Monin (2010)

%==============================================================================
% MA-EPI-35: Identifiable Victim Effect
%==============================================================================

\subsubsection*{MA-EPI-35: Identifiable Victim Effect}

\textbf{Statement:} Ein identifizierbares Opfer erzeugt mehr Empathie als statistische Opfer.

\textbf{Formel:}
\begin{equation}
U_{\text{empathy}}(1_{\text{identified}}) > U_{\text{empathy}}(N_{\text{statistical}})
\end{equation}

\textbf{Beschreibung:}
``Ein Toter ist eine Tragödie, eine Million Tote sind eine Statistik'' (Stalin zugeschrieben):
\begin{itemize}
    \item Spendenbereitschaft höher für Einzelschicksale
    \item Namen und Gesichter aktivieren Empathie
    \item Abstrakte Zahlen lassen kalt
\end{itemize}

\textbf{BCM-Relevanz:}
Erklärt asymmetrische Ressourcenallokation. Medial präsente Einzelfälle bekommen mehr als statistisch größere Probleme.

\textbf{BEATRIX-Relevanz:}
BEATRIX nutzt Identifiable Victim Effect:
\begin{itemize}
    \item Personalisierte Beispiele statt Statistiken
    \item ``Das könnte [Name] sein'' statt ``X\% Betroffene''
    \item Aber: Ethische Grenzen der Emotionalisierung
\end{itemize}

\textbf{Konsequenz bei Fehlen:}
Empathie wäre proportional zu Anzahl Betroffener. Fundraising-Strategien wären anders.

\textbf{Validiert:} KNU, IDN

\textbf{Abhängigkeiten:} MA-EPI-06

\textbf{Literatur:} Small \& Loewenstein (2003); Slovic (2007)

%==============================================================================
% MA-EPI-36: Scope Insensitivity
%==============================================================================

\subsubsection*{MA-EPI-36: Scope Insensitivity}

\textbf{Statement:} Bewertungen skalieren nicht proportional mit der Anzahl Betroffener.

\textbf{Formel:}
\begin{equation}
\frac{\partial U}{\partial N} \rightarrow 0 \quad \text{für } N \rightarrow \infty
\end{equation}

\textbf{Beschreibung:}
Die Zahlungsbereitschaft für Rettung von 2000 Vögeln ist kaum höher als für 200:
\begin{itemize}
    \item WTP für 2k, 20k, 200k Vögel: \$80, \$78, \$88
    \item ``Psychic numbing'': Große Zahlen werden nicht gefühlt
    \item 1 \approx 2 \approx viele
\end{itemize}

\textbf{BCM-Relevanz:}
Kollektiver Nutzen (KNU) wird systematisch unterschätzt. Je größer das Problem, desto mehr wird ignoriert (Klimawandel, Pandemie).

\textbf{BEATRIX-Relevanz:}
BEATRIX korrigiert für Scope Insensitivity:
\begin{itemize}
    \item Große Zahlen herunterbrechen (``pro Kopf'')
    \item Visualisierungen statt abstrakte Zahlen
    \item KNU explizit hochrechnen
\end{itemize}

\textbf{Konsequenz bei Fehlen:}
Nutzen wäre linear in Anzahl Betroffener. Kollektive Probleme würden angemessen gewichtet.

\textbf{Validiert:} KNU

\textbf{Abhängigkeiten:} MA-EPI-35

\textbf{Literatur:} Kahneman (1986); Desvousges et al. (1993)

%==============================================================================
% MA-EPI-37: Affect Heuristic
%==============================================================================

\subsubsection*{MA-EPI-37: Affect Heuristic}

\textbf{Statement:} Emotionale Reaktionen bestimmen Risiko- und Nutzenurteile.

\textbf{Formel:}
\begin{equation}
\hat{U}_{\text{risk}} = f(\text{Affect}), \quad \hat{U}_{\text{benefit}} = g(\text{Affect})
\end{equation}
\begin{equation}
\text{Corr}(\hat{U}_{\text{risk}}, \hat{U}_{\text{benefit}}) < 0
\end{equation}

\textbf{Beschreibung:}
Die Affect Heuristic ersetzt analytische Bewertung durch Bauchgefühl:
\begin{itemize}
    \item Positives Gefühl → niedriges Risiko, hoher Nutzen
    \item Negatives Gefühl → hohes Risiko, niedriger Nutzen
    \item Dies erzeugt negative Korrelation, die objektiv nicht existiert
\end{itemize}

\textbf{BCM-Relevanz:}
Erklärt emotionale Bewertung von Technologien (Kernenergie: emotional negativ → Risiko hoch, Nutzen niedrig; Solar: umgekehrt). Emotions-basierte Interventionen nutzen dies.

\textbf{BEATRIX-Relevanz:}
BEATRIX modelliert Affect:
\begin{itemize}
    \item Welche Emotionen sind mit dem Verhalten assoziiert?
    \item Können positive Affekte aktiviert werden?
    \item Emotionale vs. analytische Persona-Typen
\end{itemize}

\textbf{Konsequenz bei Fehlen:}
Risiko und Nutzen wären unabhängig bewertet. Emotionale Komponente von Entscheidungen unsichtbar.

\textbf{Validiert:} AWX, INU

\textbf{Abhängigkeiten:} MA-EPI-05

\textbf{Literatur:} Slovic et al. (2002); Finucane et al. (2000)

\bigskip
\noindent\rule{\textwidth}{0.4pt}
\textit{Ende der EPI-Axiome. Alle 37 Epistemologie-Axiome sind nun vollständig dokumentiert.}
